\chapter{Lean II: definitions, dependencies, and trust}\label{ch:lean2}
\goal{Use Lean to bring out the choices hidden in a mathematical statement. Before trusting an accepted proof, read the theorem's hypotheses and the definitions it uses, and find out whether the proof depends on unfinished steps. Learn to repair a failing proof without losing the statement you meant to prove.}

\section{The statement must be right before the proof can help}\label{ch16:sec:statement}
When Lean accepts a proof without any error or warning, it guarantees one thing: the conclusion, exactly as written in Lean, follows from the hypotheses, exactly as written in Lean. It does not guarantee that the Lean statement says what anyone meant. That statement translates a mathematical sentence, from a paper, an exercise, or a request to an AI assistant, and Lean never sees the original, so it cannot notice a mistranslation. Comparing the formal statement with the intended one is the reader's job, not the proof assistant's~\cite{leanvalidate}.

For example, Section~\ref{ch07:sec:seymour} defined the second out-neighborhood $N_2^+(v)$ as the set of vertices at directed distance exactly two from $v$. A tempting shortcut is to take all the endpoints $w$ of directed walks $v\to u\to w$. Lean would accept correct proofs about either description, but they describe different sets. In the graph of that section, with edges $a\to b$, $b\to c$, $a\to c$, and $c\to d$, the two-step walks from $a$ end at $c$ and $d$, but the edge $a\to c$ puts $c$ at distance one, so $N_2^+(a)=\{d\}$. Since $|N_1^+(a)|=2$, the vertex $a$ fails the inequality $|N_2^+(v)|\ge|N_1^+(v)|$ of Seymour's conjecture but passes it with the shortcut set $\{c,d\}$. In general the shortcut set contains $N_2^+(v)$, so it can only make the inequality easier to satisfy, and a ``proof of Seymour's conjecture'' that used it would prove a weaker statement, as that section warned.

So before reading any tactics, compare the two statements in a table: in one column, list the objects, hypotheses, and conclusion of the mathematical statement; in the other, write what represents each of them in Lean. This is the statement contract of Habit~2 (Section~\ref{sec:magnitude}), written out for a formal statement; Section~\ref{ch16:sec:surj} gives a complete example. While you fill it in, check the places where translations most often go wrong; most can be seen in Seymour's conjecture: \emph{every nonempty finite oriented graph has a vertex $v$ with $|N_2^+(v)|\ge|N_1^+(v)|$}.
\begin{itemize}
\item \emph{Domains.} Does each variable range over the intended set? In Section~\ref{ch15:sec:exists}, the expression $7-10$ meant $0$ for natural numbers and $-3$ for integers.
\item \emph{Strict and non-strict inequalities.} With $>$ in place of $\ge$, the conjecture would be false: a graph with one vertex and no edges is a nonempty finite oriented graph, and its only vertex has $|N_2^+(v)|=0=|N_1^+(v)|$.
\item \emph{Nonemptiness.} A graph with no vertices has no vertex with the required property, so without the word ``nonempty'' the conjecture would be false.
\item \emph{Excluded cases.} The conditions $w\ne v$ and $w\notin N_1^+(v)$ in the definition of $N_2^+(v)$ are exactly what the shortcut leaves out.
\item \emph{The order of quantifiers.} The conjecture claims that \emph{some} vertex satisfies the inequality, not that every vertex does: in the graph above, $a$ fails it, while the vertex $d$, which has no out-neighbors, satisfies it with $0\ge0$. Likewise, ``for every $y\in B$ there is an $x\in A$ with $f(x)=y$'' says that $f$ is surjective, whereas ``there is an $x\in A$ with $f(x)=y$ for every $y\in B$'' asks a single input to be sent to every element of $B$ at once, which is impossible once $B$ has two elements (Section~\ref{ch02:sec:quantifiers}).
\end{itemize}

\subsection*{Finding the statement and its definitions}
To fill in the table for an unfamiliar Lean theorem, read its parameters and the claim after its colon (Section~\ref{ch15:sec:symbols}). Then look up each definition that the statement uses and read its \emph{body}, the part after its own \code{:=}. Only then turn to the proof. The name of a definition says no more than the name \code{all\_problems\_solved} of a theorem whose claim is only \code{0 = 0}: a definition called \code{Second} could describe the second out-neighborhood with or without the conditions $w\ne v$ and $w\notin N_1^+(v)$, and only its body tells you which.

\section{A complete existence proof in Lean}\label{ch16:sec:surj}
In Section~\ref{ch15:sec:exists}, the witness for an existence statement was easy to spot. In many existence proofs the witness has to be built from other information. Our example is Exercise~\ref{ex:3.2}: if $f:A\to B$ and $g:B\to C$ are surjective, then so is $g\circ f$. Recall from Section~\ref{ch03:sec:injsurj} that $f$ is surjective when every target has a preimage:
\[
\forall b\in B\ \exists a\in A:\qquad f(a)=b.
\]
Here is this definition in Lean, followed by the theorem and its proof.
\par\noindent\begin{minipage}{\linewidth}
\begin{lstlisting}
def Surj {A B : Type} (f : A -> B) : Prop :=
  forall y : B, Exists (fun x : A => f x = y)

theorem surj_comp {A B C : Type}
    (f : A -> B) (g : B -> C)
    (hf : Surj f) (hg : Surj g) :
    Surj (fun x => g (f x)) := by
  intro c
  cases hg c with
  | intro b hb =>
    cases hf b with
    | intro a ha =>
      refine Exists.intro a ?_
      show g (f a) = c
      rw [ha, hb]
\end{lstlisting}
\end{minipage}\par

\subsection*{The statement}
As Section~\ref{ch16:sec:statement} recommends, we read the statement before the proof. The body of the definition says that \code{Surj f} is the proposition \code{forall y : B, Exists (fun x : A => f x = y)}: for every \code{y} in \code{B}, some \code{x} in \code{A} satisfies \code{f x = y}. The body shows the domain \code{A}, the codomain \code{B}, and the order of the quantifiers. The target \code{y} is chosen first, and the input \code{x} is found afterward, so it may depend on \code{y}. The table compares the theorem with the mathematical statement.
\begin{center}
\begin{tabularx}{\linewidth}{@{}>{\raggedright\arraybackslash}X>{\raggedright\arraybackslash}X@{}}
\toprule
Mathematics & Lean\\
\midrule
sets $A$, $B$, $C$ & types \code{A B C : Type}, as implicit parameters\\
functions $f:A\to B$ and $g:B\to C$ & \code{f : A -> B} and \code{g : B -> C}\\
hypothesis: $f$ is surjective & \code{hf : Surj f}\\
hypothesis: $g$ is surjective & \code{hg : Surj g}\\
the composition $g\circ f$ & \code{fun x => g (f x)}\\
conclusion: $g\circ f$ is surjective & \code{Surj (fun x => g (f x))}\\
\bottomrule
\end{tabularx}
\end{center}
Every part of the mathematical statement has a partner in Lean, and nothing extra has been assumed. Section~\ref{ch16:sec:irrelevant} shows what an extra assumption can do.

\subsection*{The proof, line by line}
Because \code{Surj} is defined as a ``for all'' statement, the tactic \code{intro} can look inside the definition, just as it looked inside \code{Inj} in Section~\ref{ch15:sec:injcomp}. The line \code{intro c} introduces an arbitrary target \code{c} in \code{C}, and the infoview shows the goal
\begin{center}
$\exists$\,\code{x, (fun x => g (f x)) x = c}
\end{center}
As at the end of Section~\ref{ch15:sec:injcomp}, Lean shows the composition applied to \code{x} without evaluating it. In words, the goal asks for some $x$ in $A$ with $g(f(x))=c$. We do not yet know which element to use. Our only sources of elements are the two hypotheses, and we start with the one about $g$, because the outputs of $g$ lie in $C$, where the target $c$ lives.

Since \code{hg} proves a ``for all'' statement, it can be applied to \code{c}, just as the hypothesis \code{hg : Inj g} was applied to elements of \code{B} in Section~\ref{ch15:sec:injcomp}. The result \code{hg c} proves that some element of \code{B} is sent to \code{c} by \code{g}. Like all existence evidence, it consists of a witness together with a proof that the witness has the property. The lines
\par\noindent\begin{minipage}{\linewidth}
\begin{lstlisting}
  cases hg c with
  | intro b hb =>
\end{lstlisting}
\end{minipage}\par
use the tactic \code{cases}, which takes evidence apart. Here they mean: unpack this existence evidence, call the witness \code{b} and the proof of its property \code{hb}, and continue with the indented lines below. These indented lines form a \emph{branch} of the proof. Inside this branch the facts \code{b : B} and \code{hb : g b = c} are available, and the goal has not changed. The word \code{intro} after the vertical bar is not the tactic \code{intro}. It refers to \code{Exists.intro}, the only way to build existence evidence, which Section~\ref{ch15:sec:exists} used to combine a witness with a proof. Writing \code{| intro b hb} says: this evidence was built as \code{Exists.intro b hb}, so take it apart into those two pieces. The infoview labels the goal inside this branch \code{case intro}, after the same word; inside the second unpacking, described next, the label becomes \code{case intro.intro}.

The second unpacking does the same with \code{hf b}. Since $b$ lies in $B$, the surjectivity of $f$ applies to it, and it supplies a witness \code{a : A} together with \code{ha : f a = b}. Inside this second branch the available information is
\[
 b\in B,\quad g(b)=c,\qquad a\in A,\quad f(a)=b.
\]
Now we have a candidate input, $a$. The line \code{refine Exists.intro a ?\_} proposes $a$ as the witness. The tactic \code{refine} works like \code{exact}, except that its argument may contain the placeholder \code{?\_} for a part of the proof that has not been supplied yet; that part becomes the new goal. Here the missing part is the proof that $a$ has the required property, and Lean displays the new goal as
\par\noindent\begin{minipage}{\linewidth}
\begin{lstlisting}
(fun x => g (f x)) a = c
\end{lstlisting}
\end{minipage}\par
Again the composition is shown without being evaluated: this is the rule $x\mapsto g(f(x))$, applied to the input $a$. The line \code{show g (f a) = c} replaces the goal by its beta-reduced form. The tactic \code{show} accepts a new form of the goal only when Lean can check, by reducing expressions, that it is the same proposition as the old one. Here the check succeeds, because beta-reducing \code{(fun x => g (f x)) a} gives \code{g (f a)}.

The \code{show} line is needed. The next tactic, \code{rw [ha, hb]}, rewrites by searching the goal for the written expression \code{f a}, the left side of \code{ha}. In \code{(fun x => g (f x)) a} that expression does not occur: inside the rule, the goal contains \code{f x}, with the variable \code{x}. Without the \code{show} line, Lean stops at the \code{rw} line with the error
\par\noindent\begin{minipage}{\linewidth}
\begin{lstlisting}
tactic 'rewrite' failed, did not find instance of the pattern in the target expression
  f a
\end{lstlisting}
\end{minipage}\par
After \code{show}, the expression \code{f a} is visible. Rewriting with \code{ha : f a = b} turns the goal into \code{g b = c}. Rewriting with \code{hb : g b = c} then turns it into \code{c = c}, which \code{rw} closes by itself (Section~\ref{ch15:sec:equality}).

The table records the proof state after each step.
\begin{center}
\begin{tabularx}{\linewidth}{@{}>{\raggedright\arraybackslash}p{0.27\linewidth}>{\raggedright\arraybackslash}X>{\raggedright\arraybackslash}X@{}}\toprule
Step & Facts added & Goal afterwards\\\midrule
Start & \code{hf : Surj f} and \code{hg : Surj g} & \code{Surj (fun x => g (f x))}\\
\code{intro c} & \code{c : C} & $\exists$\,\code{x, (fun x => g (f x)) x = c}\\
unpack \code{hg c} & \code{b : B} and \code{hb : g b = c} & unchanged\\
unpack \code{hf b} & \code{a : A} and \code{ha : f a = b} & unchanged\\
\code{refine Exists.intro a ?\_} & none & \code{(fun x => g (f x)) a = c}\\
\code{show g (f a) = c} & none & \code{g (f a) = c}\\
\code{rw [ha, hb]} & none & none: the proof is complete\\\bottomrule
\end{tabularx}
\end{center}

The witnesses are found in the opposite order to the one in which the functions are applied. The composition applies $f$ first and $g$ second. The proof first finds an input to $g$ that reaches the final target $c$, and then an input to $f$ that reaches the intermediate value $b$. The solution of Exercise~\ref{ex:3.2} finds them in the same order.

\pauseq{Could the proof begin by unpacking \code{hf} instead of \code{hg}?}{No. The hypothesis \code{hf} says that every element of $B$ has a preimage under $f$, so it can only be applied to an element of \code{B}. At the start, the only element available is the target \code{c}, which lies in \code{C}. Lean rejects \code{hf c} with an ``application type mismatch'' error, because \code{c} has type \code{C} where an element of \code{B} is expected. An element of $B$ becomes available only after \code{hg c} has been unpacked and has supplied \code{b}. The types of the hypotheses force the order of the steps.}

\subsection*{The ordinary proof beside the Lean proof}
Here is the same argument written on paper.
\begin{quote}
Let $c\in C$. Since $g$ is surjective, there is $b\in B$ with $g(b)=c$. Since $f$ is surjective, there is $a\in A$ with $f(a)=b$. Then $(g\circ f)(a)=g(f(a))=g(b)=c$. Since $c$ was arbitrary, $g\circ f$ is surjective.
\end{quote}
Each sentence has a partner in the Lean proof. ``Let $c\in C$'' is \code{intro c}. The two sentences that begin with ``Since $g$'' and ``Since $f$'' are the two unpackings; each uses a hypothesis to obtain a named witness. The chain of equalities proposes $a$ as the witness and checks it, which Lean does with \code{refine}, \code{show}, and \code{rw}. The last sentence has no separate line in Lean. Because \code{c} was introduced as an arbitrary element, Lean already knows that the argument covers every \code{c}.

The indentation shows where each name may be used. The names \code{b} and \code{hb} exist only inside the branch that follows \code{| intro b hb =>}, and \code{a} and \code{ha} exist only inside the branch nested within it. Before the first unpacking, the proof knows that some element of $B$ is sent to $c$, but it has no name for that element and cannot use it. A written proof behaves the same way: the sentence ``there is $b\in B$ with $g(b)=c$'' introduces the name $b$ and records the one fact that is known about it. Notice also that $b$ need not be the only element of $B$ with this property, since $c$ may have several preimages under $g$. The proof uses nothing about $b$ except $g(b)=c$, so it works whichever of them $b$ stands for. This is the care that Section~\ref{ch04:sec:unique} asked for when it distinguished ``suppose $x$ is a solution'' from ``let $x$ be the solution.''

\section{Sets as membership predicates}\label{ch16:sec:predicates}
A \emph{predicate} on a set $A$ is a rule that turns each element of $A$ into a proposition. For example, ``$n$ is even'' is a predicate on the integers, giving statements such as ``$6$ is even'' (true) and ``$7$ is even'' (false). In Lean, a predicate on a type \code{A} is a function of type \code{A -> Prop}; you met one in Section~\ref{ch15:sec:exists}, where the input to \code{Exists} was \code{fun n : Nat => n + 1 = 4}.

A predicate $P$ on $A$ describes the set $\{x\in A:P(x)\}$ of Section~\ref{ch03:sec:sets}, and conversely a set $S\subseteq A$ gives the predicate ``$x\in S$.'' So a subset of a type \code{A} can be represented in Lean by a predicate \code{s : A -> Prop}, where \code{s x} means that $x$ belongs to the set; nothing is lost, because for every $x$ the question ``does $x$ belong?'' has the same answer in both descriptions. For example, the set of even natural numbers is represented by \code{fun n : Nat => Exists (fun k : Nat => n = 2 * k)}.

Here are intersection and preimage, defined by their membership conditions, with a theorem that connects them:
\par\noindent\begin{minipage}{\linewidth}
\begin{lstlisting}
def meet {A : Type} (s t : A -> Prop) : A -> Prop :=
  fun x => And (s x) (t x)

def preimage {A B : Type} (f : A -> B)
    (s : B -> Prop) : A -> Prop :=
  fun x => s (f x)

theorem preimage_meet {A B : Type} (f : A -> B)
    (s t : B -> Prop) (x : A) :
    Iff (preimage f (meet s t) x)
      (meet (preimage f s) (preimage f t) x) := by
  rfl
\end{lstlisting}
\end{minipage}\par

If \code{s} and \code{t} describe the sets $S$ and $T$, then \code{meet s t} holds at \code{x} when both \code{s x} and \code{t x} hold, so it describes $S\cap T$; and \code{preimage f s} holds at \code{x} when $f(x)\in S$, so it describes the preimage $f^{-1}(S)$. Since \code{Iff P Q} is Lean's way of writing $P\Leftrightarrow Q$, the theorem \code{preimage\_meet} is the identity
\[
f^{-1}(S\cap T)=f^{-1}(S)\cap f^{-1}(T)
\]
from Section~\ref{ch03:sec:images}, stated one element at a time: for each \code{x}, membership in the left side is equivalent to membership in the right side, which is what it means for the two sets to be equal.

Why does \code{rfl} prove it? The left side \code{preimage f (meet s t) x} unfolds to \code{meet s t (f x)} and then to \code{And (s (f x)) (t (f x))}. The right side \code{meet (preimage f s) (preimage f t) x} unfolds to \code{And (preimage f s x) (preimage f t x)} and then to the very same \code{And (s (f x)) (t (f x))}. Just as \code{rfl} proves an equality whose two sides reduce to the same expression (Section~\ref{ch15:sec:symbols}), it proves an equivalence whose two sides reduce to the same proposition, unfolding the definitions by itself. These unfoldings are exactly the steps ``by definition'' of the chain of equivalences in Section~\ref{ch03:sec:images}.

\pauseq{Suppose the right side of \code{preimage\_meet} were written \code{meet (preimage f t) (preimage f s) x}, with \code{s} and \code{t} exchanged. The statement is still true. Would \code{rfl} still prove it?}{No. The right side now unfolds to \code{And (t (f x)) (s (f x))}, which is equivalent to \code{And (s (f x)) (t (f x))} but not the same expression, so \code{rfl} fails, with an error saying that the two sides are ``not definitionally equal.'' A proof would have to swap the two halves of the conjunction in each direction of the equivalence, for instance with the theorem \code{and\_swap} of Section~\ref{ch15:sec:conjunction}.}

Lean writes the disjunction $P\lor Q$ of Section~\ref{ch02:sec:connectives} as \code{Or P Q}. Had \code{meet} been defined with \code{Or} in place of \code{And}, the theorem \code{preimage\_meet} would keep its name and its proof \code{rfl}, and Lean would accept it; but it would be the statement about unions from Exercise~\ref{ex:16.2}. The bodies of the definitions, not their names, decide what a statement says.

\section{Encoding the exclusions in a second out-neighborhood}\label{ch16:sec:second}
A directed graph whose vertices form a type \code{V} can be described in Lean by its edge relation \code{r : V -> V -> Prop}. As in Section~\ref{ch15:sec:injcomp}, a type with two arrows describes functions of two inputs: \code{r} takes two vertices \code{v} and \code{w} and returns the proposition \code{r v w}, which we read as $v\to w$. A relation of this kind can describe more than the directed graphs of Section~\ref{ch07:sec:sink}, because nothing in it rules out a loop $v\to v$, which Definition~\ref{ch07:def:directed} does not allow. Like a directed graph, it also allows a pair of opposite edges $v\to w$ and $w\to v$, which an oriented graph does not. When a statement needs the graph to have no loops, or to be oriented, it must say so in extra hypotheses.

Recall from Section~\ref{ch07:sec:seymour} that $w$ belongs to the second out-neighborhood $N_2^+(v)$ when three conditions hold: $w\ne v$; $w$ is not an out-neighbor of $v$; and $v\to u\to w$ for some vertex $u$. In Lean, the three conditions become a definition that turns each pair of vertices $v$, $w$ into a proposition:
\par\noindent\begin{minipage}{\linewidth}
\begin{lstlisting}
def Second {V : Type} (r : V -> V -> Prop)
    (v w : V) : Prop :=
  And (Not (w = v)) (And (Not (r v w))
    (Exists (fun u : V => And (r v u) (r u w))))
\end{lstlisting}
\end{minipage}\par
Here \code{Not P} is the negation of \code{P}. So \code{Not (w = v)} expresses $w\ne v$, and \code{Not (r v w)} says that there is no edge $v\to w$. The last part says that some vertex \code{u} satisfies both \code{r v u} and \code{r u w}. It asks for at least one route $v\to u\to w$, not for a unique one; if several vertices $u$ give a route, $w$ is still a single member of $N_2^+(v)$ (Section~\ref{ch07:sec:seymour}).

Lean also has the usual symbols for these connectives, and the infoview uses them. When the definition is unfolded, for instance with the tactic \code{unfold Second}, the infoview displays the body of \code{Second r v w} as
\begin{center}
$\neg$\,\code{w = v} $\land$ $\neg$\,\code{r v w} $\land$ $\exists$\,\code{u, r v u} $\land$ \code{r u w}
\end{center}
with no parentheses at all. Three conventions tell you how to read it. First, $\neg$ applies only to the proposition immediately after it, such as the whole equation \code{w = v}, and it stops at the next $\land$. So $\neg$\,\code{w = v} $\land$ $\neg$\,\code{r v w} means $(\neg$\,\code{w = v}$)\land(\neg$\,\code{r v w}$)$, which says that $w\ne v$ and that there is no edge $v\to w$. Second, $\exists$\,\code{u,} reaches as far to the right as possible, so it covers both \code{r v u} and \code{r u w}, just as \code{Exists} does in the definition. Third, $\land$ groups to the right, so that $P\land Q\land R$ means $P\land(Q\land R)$. In the Lean code of this book we write the names \code{And}, \code{Not}, and \code{Exists} instead, because they are easy to type and, together with the parentheses, they make the grouping visible.

To read the nesting, give the three conditions temporary names:
\[
 P:\ w\ne v,\qquad Q:\ \text{there is no edge }v\to w,\qquad
 R:\ v\to u\to w\text{ for some vertex }u.
\]
The definition has the shape \code{And P (And Q R)}. A proof \code{h} of it is a pair (Section~\ref{ch15:sec:conjunction}). Its first component is a proof of $P$, and its second component is itself a pair, made of a proof of $Q$ and a proof of $R$. So \code{h.left} is a proof of $P$, and \code{h.right} is a proof of \code{And Q R}. Taking components once more, \code{h.right.left} is a proof of $Q$, and \code{h.right.right} is a proof of $R$. Taking a component of a pair in this way is called a \emph{projection}.

\pauseq{Take the oriented graph of Section~\ref{ch07:sec:seymour} with edges $a\to b$, $b\to c$, $a\to c$, and $c\to d$, and let $v=a$. For each vertex $w$, decide which of the conditions $P$, $Q$, $R$ hold. For which $w$ does \code{Second r a w} hold?}{For $w=a$, condition $P$ fails. Condition $Q$ holds, because there is no edge $a\to a$, but $R$ fails: the out-neighbors of $a$ are $b$ and $c$, and neither has an edge to $a$. For $w=b$, condition $P$ holds, but $Q$ fails because of the edge $a\to b$. Condition $R$ fails as well, since neither $b$ nor $c$ has an edge to $b$. For $w=c$, condition $P$ holds and so does $R$, through the route $a\to b\to c$, but $Q$ fails because of the edge $a\to c$. For $w=d$, all three hold: $d\ne a$, there is no edge $a\to d$, and $a\to c\to d$ is a route. So \code{Second r a w} holds only for $w=d$, in agreement with $N_2^+(a)=\{d\}$. The vertex $c$ is excluded by condition $Q$ alone, one of the two conditions that the shortcut of Section~\ref{ch16:sec:statement} leaves out.}

A small but useful consequence of the definition is that a second out-neighbor is never a first out-neighbor:
\par\noindent\begin{minipage}{\linewidth}
\begin{lstlisting}
theorem second_not_first {V : Type}
    (r : V -> V -> Prop) (v w : V)
    (h : Second r v w) : Not (r v w) := by
  exact h.right.left
\end{lstlisting}
\end{minipage}\par
The proof is the projection \code{h.right.left} described above: it extracts the proof of $Q$, which is exactly the claim \code{Not (r v w)}. If the definition grouped the conditions differently, the projections would change with it. Exercise~\ref{ex:16.3} asks you to extract the other two conditions.

The formal statement of Seymour's conjecture in the repository of Section~\ref{ch07:sec:seymour} needs more than this one predicate. Like us, it uses a type \code{V} of vertices and an edge relation \code{r}. It adds hypotheses saying that there are finitely many vertices, that there is at least one, and that \code{r} is oriented: no vertex has an edge to itself, and no two vertices are joined in both directions. It also needs a way to count the members of each out-neighborhood. Its condition for second out-neighbors, however, is exactly the one above: $w\ne v$, there is no edge $v\to w$, and $v\to u\to w$ for some $u$~\cite{oai173}. The repository's statement file~\cite{oaistatement}, the Lean file named in its configuration~\cite{oaiconfig}, even joins the three conditions by $\land$ in the same order. So the piece of the statement that the shortcut gets wrong can be checked on its own, before anyone studies the long proof. Section~\ref{ch17:sec:research} reads the full formal statement piece by piece.

\section{A successful proof can still be irrelevant}\label{ch16:sec:irrelevant}
Lean accepts the following theorem without any error or warning:
\par\noindent\begin{minipage}{\linewidth}
\begin{lstlisting}
theorem impossible_context (h : False) : (0 : Nat) = 1 := by
  exact False.elim h
\end{lstlisting}
\end{minipage}\par

Its conclusion, $0=1$, is false. Yet the theorem is correct, because the conclusion is not all that it says. The parameter \code{h : False} is a hypothesis, and \code{False} is the proposition that is never true, so it has no proof. The theorem says: \emph{if} we had a proof of \code{False}, then $0=1$ would follow. This implication is true for the reason that Section~\ref{ch02:sec:contrapositive} gave for vacuous truth: its hypothesis can never be satisfied. The proof uses the rule \code{False.elim}, which turns a proof of \code{False} into a proof of any proposition whatsoever. Since nobody can ever supply the hypothesis \code{h}, the theorem can never be used to conclude that $0=1$. A reader who looked only at the conclusion would completely misunderstand what was proved.

The same thing can happen less visibly when an assistant ``repairs'' a proof that Lean rejects by changing the statement, most often by adding a hypothesis. The new statement may well be true and its proof correct, and still not be the statement that was asked for. Here are two examples.
\begin{itemize}
\item The claim ``for all real numbers $x$ and $y$, if $x^2=y^2$, then $x=y$'' is false: take $x=1$ and $y=-1$. Adding the hypotheses $x\ge0$ and $y\ge0$ makes it true, because the square map is injective on $[0,\infty)$ (Section~\ref{ch03:sec:injsurj}). The repaired claim is a correct theorem, but it says nothing about negative numbers. Whether the repair is acceptable depends on whether those hypotheses were part of the intended statement.
\item An existence theorem can be made trivial by assuming the very thing to be proved. Here is \code{surj\_comp} with one extra hypothesis:
\par\noindent\begin{minipage}{\linewidth}
\begin{lstlisting}
theorem surj_comp_assumed {A B C : Type}
    (f : A -> B) (g : B -> C)
    (hf : Surj f) (hg : Surj g)
    (h : Surj (fun x => g (f x))) :
    Surj (fun x => g (f x)) := by
  exact h
\end{lstlisting}
\end{minipage}\par
The new hypothesis \code{h} is the conclusion itself, so the proof only has to repeat it. Lean accepts this theorem. It does print a warning that the variables \code{hf} and \code{hg} are unused, a useful hint that the proof never needed the real hypotheses. But a warning is not an error, and a changed statement does not always leave such a trace: \code{impossible\_context} produces no warning at all.
\end{itemize}
In both examples the repaired statement is true and has a correct proof, but the mathematical task has changed. So when a proof has been repaired, first compare the new statement with the old one, hypothesis by hypothesis and then conclusion with conclusion, as in the table of Section~\ref{ch16:sec:surj}. Only then look at how the proof itself changed.

\section{Placeholders and axiom dependencies}\label{ch16:sec:axioms}
While a proof is being developed, a gap can be filled temporarily with the word \code{sorry}. Lean then accepts the declaration with the warning \code{declaration uses 'sorry'}, treating the missing step as if it had been assumed; the proof is not complete. In fact \code{sorry} closes any goal, even a false one: Lean accepts \code{theorem bad : (0 : Nat) = 1 := by sorry}, with the same warning.

The warning appears only at the declaration that contains \code{sorry}. If the last line of \code{surj\_comp} is replaced by \code{sorry}, Lean warns at \code{surj\_comp}, but a later theorem whose proof applies \code{surj\_comp} gets no warning of its own, although it rests on the same gap. To find out whether a theorem depends on a gap anywhere, directly or through other theorems, use the command
\par\noindent\begin{minipage}{\linewidth}
\begin{lstlisting}
#print axioms theoremName
\end{lstlisting}
\end{minipage}\par
with the theorem's name in place of \code{theoremName}. It lists every \emph{axiom}, that is, every statement accepted without proof, on which the theorem's proof ultimately relies. A \code{sorry} anywhere among its dependencies shows up as the axiom \code{sorryAx}.

Not every axiom signals a problem. Lean's foundations include three standard axioms that ordinary mathematics uses all the time:
\begin{itemize}
\item \code{propext}, short for \emph{propositional extensionality}: two propositions that are logically equivalent are equal, so that if $P\Leftrightarrow Q$, then $P$ may be replaced by $Q$ anywhere;
\item \code{Quot.sound}: in a quotient, the set of equivalence classes of Section~\ref{ch06:sec:residues}, two related elements have the same class, as in $[1]_4=[5]_4$, since $4$ divides $5-1$;
\item \code{Classical.choice}: once it has been proved that an object of some kind exists, Lean may pick one, even when no particular one can be described. From it, with the other two axioms, Lean derives the principle that every proposition is either true or false, which proof by contradiction uses in its last step (Section~\ref{ch04:sec:contrapositive}); so a proof by contradiction about an arbitrary proposition usually depends on \code{Classical.choice}.
\end{itemize}
Seeing these in the list is normal. What matters is to know what each listed axiom says~\cite{leanaxioms}, and to notice one that does not belong, such as \code{sorryAx}.

For three of the examples in Chapters~\ref{ch:lean1} and~\ref{ch:lean2}, the commands
\par\noindent\begin{minipage}{\linewidth}
\begin{lstlisting}
#print axioms MathCourse.inj_comp
#print axioms MathCourse.surj_comp
#print axioms MathCourse.preimage_meet
\end{lstlisting}
\end{minipage}\par
print the three lines
\par\noindent\begin{minipage}{\linewidth}
\begin{lstlisting}
'MathCourse.inj_comp' does not depend on any axioms
'MathCourse.surj_comp' does not depend on any axioms
'MathCourse.preimage_meet' does not depend on any axioms
\end{lstlisting}
\end{minipage}\par
So these proofs use only the rules of Lean's logic and our own definitions. (Appendix~\ref{appD:sec:lean} explains where such commands go in the file.)

An axiom list is no substitute for reading the statement. The theorems \code{impossible\_context} and \code{surj\_comp\_assumed} of Section~\ref{ch16:sec:irrelevant} depend on no axioms at all, yet the first assumes \code{False} and the second its own conclusion. To trust a theorem, you need to know both the assumptions in its statement and the axioms its proof depends on.

When you check a file, read everything Lean reports, warnings included, and the output of any \code{\#print axioms} commands; and wait until Lean has finished the whole file, since a line it has not yet reached shows no error. Proofs from sources you do not trust call for the stronger procedures in the official documentation~\cite{leanvalidate}; Section~\ref{ch17:sec:research} describes the check used for Seymour's conjecture.

\subsection*{How far a check reaches}
A different question is how many cases a check covers, even when it finishes without any problem. Compare three kinds of check.
\begin{itemize}
\item A Python program, such as the one in Section~\ref{appD:sec:experiments}, can run through all $2^4=16$ sign sequences of length $4$ and report that exactly $8$ of them have $C_1=C_2=C_3=0$ (Section~\ref{ch14:sec:experiment}). This covers $16$ cases and says nothing about other lengths, and no checker compares the code with the definition of $C_s$.
\item A Lean proof about specific objects, such as \code{twice\_two} in Section~\ref{ch15:sec:symbols}, is checked step by step, but it covers only the case it states: \code{twice 2 = 4}, and nothing about \code{twice 3}.
\item A Lean proof of a universal statement, such as \code{surj\_comp}, covers every case at once: all types $A$, $B$, $C$ and all surjective functions between them, infinitely many cases that no program could test.
\end{itemize}
In all three, the result concerns the statement as written: code or a definition that encodes the wrong quantity, through a slip in the wrap-around or by using the aperiodic correlation where the periodic one was meant, gives a correct answer to the wrong question.

\section{A debugging sequence that does not lose the mathematics}\label{ch16:sec:debugging}
When Lean rejects a line, it is tempting to try one tactic after another, or to ask an assistant for anything that Lean will accept. Either habit can end with a file that Lean accepts and a statement that nobody wanted. Finding and removing the cause of an error is called \emph{debugging}; the following sequence keeps the mathematics in view while you do it.
\begin{enumerate}
\item \emph{Keep the statement fixed.} Do not edit the theorem's statement while you look for the problem.
\item \emph{Read the proof state.} Look at the goal, the available facts, and the error message.
\item \emph{Translate into ordinary mathematics.} Say in words what the goal claims and what each fact says.
\item \emph{Diagnose, then respond.} Decide which of four situations you are in.
\begin{itemize}
\item A \emph{difference of form}: the goal follows from the available facts, but something is written in a different form from the one the tactic expects. A one-line change, such as a \code{show} line, is usually enough.
\item A \emph{missing lemma}: the goal is true, but the proof needs a fact that is not yet available. Find one, read its exact statement with the command \code{\#check} (described below), and make sure that its hypotheses hold in your situation before you use it.
\item A \emph{mismatch of domains}: a variable ranges over a different set from the intended one. Go back to the statement contract and decide which domain was intended. If the statement has to change, change it on purpose and record the change, for instance in a comment above the theorem.
\item A \emph{goal that does not follow}: in some situation where all the available facts hold, the goal fails. Stop looking for tactics, since none can derive it. First ask whether an earlier step caused the problem: in Section~\ref{ch15:sec:exists}, offering the witness $2$ would have required a proof of the false equation $2+1=4$, and the remedy was a better witness. If the goal already fails to follow at the start, or after steps such as \code{intro} that only name arbitrary objects and assume the hypotheses, then the statement itself is false. Look for a \emph{counterexample}, an instance in which every hypothesis holds and the conclusion fails; what to do next is then a decision about the statement.
\end{itemize}
\end{enumerate}

\subsection*{Two worked examples}
\emph{A difference of form.} Leave out the line \code{show g (f a) = c} in the proof of \code{surj\_comp}, keeping the statement (step~1). The proof state (step~2) shows that \code{rw} did not find \code{f a} in the goal \code{(fun x => g (f x)) a = c}. Translated (step~3), this goal says $(g\circ f)(a)=c$, which follows from $f(a)=b$ and $g(b)=c$. So the diagnosis (step~4) is a difference of form, and the response is the \code{show} line, which changes the form of the goal but not its meaning (Section~\ref{ch16:sec:surj}).

\emph{A false statement.} Suppose we set out to prove \code{Surj (fun n : Nat => n + 1)}, that $n\mapsto n+1$ is surjective on the natural numbers. After \code{intro y}, the goal is
\begin{center}
$\exists$\,\code{x, (fun n => n + 1) x = y}
\end{center}
that is, for this arbitrary $y$ we must find a natural number $x$ with $x+1=y$. Suppose that every tactic we try fails. Instead of trying more, translate the goal and test a small case: for $y=0$ we would need $x+1=0$, which no natural number satisfies. Since \code{intro y} only named an arbitrary target, the problem cannot come from an earlier step: the statement itself is false, with the counterexample $y=0$, and no tactic will ever close this goal. An assistant asked only to make Lean accept the file could succeed only by changing the statement (Section~\ref{ch16:sec:irrelevant}) or by leaving a \code{sorry} (Section~\ref{ch16:sec:axioms}).

What should happen next depends on what was intended, which only a person can decide. If the claim was meant to be about natural numbers, it is false, and Lean can prove that it is false:
\par\noindent\begin{minipage}{\linewidth}
\begin{lstlisting}
theorem not_surj_succ :
    Not (Surj (fun n : Nat => n + 1)) := by
  intro h
  cases h 0 with
  | intro n hn =>
    exact Nat.succ_ne_zero n hn
\end{lstlisting}
\end{minipage}\par
In Lean, \code{Not P} is defined to mean \code{P -> False}, so \code{intro h} assumes that the map is surjective and leaves the goal \code{False}. Applying \code{h} to the target \code{0} and unpacking gives a natural number \code{n} and a proof \code{hn} that $n+1=0$; the infoview shows it unreduced, as \code{(fun n => n + 1) n = 0}, which Lean treats as the same proposition. The last step uses the library lemma \code{Nat.succ\_ne\_zero}, whose exact statement step~4 tells us to read first. To see the statement of any named theorem or lemma, put \code{\#check} followed by its name on a line of its own, outside every proof; like \code{\#eval} (Section~\ref{ch15:sec:exists}), it shows its output in the infoview. Here \code{\#check Nat.succ\_ne\_zero} prints
\begin{center}
\code{Nat.succ\_ne\_zero (n : Nat) : n.succ} $\ne$ \code{0}
\end{center}
It reads like a declaration, with the parameter \code{(n : Nat)} before the colon and the claim after it. Here \code{n.succ}, the \emph{successor} of \code{n}, is another way of writing \code{n + 1}, and \code{a} $\ne$ \code{b} means \code{Not (a = b)}. So the lemma says that $n+1\ne0$ for every natural number $n$, with no further hypotheses. Applied to our \code{n}, it proves \code{Not (n + 1 = 0)}, that is, \code{n + 1 = 0 -> False}, and applying this to \code{hn} gives the required proof of \code{False}.

If instead the claim was meant to be about integers, the domain was mistranslated: over $\ZZ$, the target $y$ is reached by the input $y-1$. Changing \code{Nat} to \code{Int} is a deliberate change of the statement, recorded as such in the comment on the first line:
\par\noindent\begin{minipage}{\linewidth}
\begin{lstlisting}
-- Statement changed on purpose: Int in place of Nat.
theorem succ_surj_int : Surj (fun n : Int => n + 1) := by
  intro y
  refine Exists.intro (y - 1) ?_
  show y - 1 + 1 = y
  exact Int.sub_add_cancel y 1
\end{lstlisting}
\end{minipage}\par
As in \code{surj\_comp}, we introduce a target, propose a witness, restate the goal in reduced form, and prove it. The library lemma \code{Int.sub\_add\_cancel a b} proves $(a-b)+b=a$ for all integers $a$ and $b$; with $a=y$ and $b=1$ it is exactly the goal. The same step would fail for natural numbers: with truncated subtraction (Section~\ref{ch15:sec:exists}), the witness for the target $0$ would be $0-1=0$, and $(0-1)+1=1\ne0$.

For these two theorems, \code{\#print axioms} lists \code{propext}, which the library lemmas \code{Nat.succ\_ne\_zero} and \code{Int.sub\_add\_cancel} use; as Section~\ref{ch16:sec:axioms} explained, this standard axiom is no cause for concern.

\section*{Exercises}
Exercises~\ref{ex:16.1}--\ref{ex:16.4} are core exercises, and Exercises~\ref{ex:16.5} and~\ref{ex:16.6} are transfer exercises. Complete solutions are in Appendix~\ref{app:solc}, under ``Week~16.''
\begin{exercise}\label{ex:16.1}In the proof of \code{surj\_comp} (Section~\ref{ch16:sec:surj}), explain what \code{hb} and \code{ha} state, and why the witness that the proof finally proposes must be an element of the domain $A$ of $f$.\end{exercise}
\begin{exercise}\label{ex:16.2}Define a predicate \code{join s t} that describes the union of the sets described by \code{s} and \code{t}, as \code{meet s t} describes their intersection. Then state and prove in Lean the corresponding fact about preimages: for each \code{x}, \code{x} belongs to the preimage of the union exactly when it belongs to the union of the preimages.\end{exercise}
\begin{exercise}\label{ex:16.3}Suppose \code{h : Second r v w}. Write a Lean proof that $w\ne v$, and a Lean proof that some vertex $u$ gives a route $v\to u\to w$.\end{exercise}
\begin{exercise}\label{ex:16.4}An assistant cannot prove a theorem, so it adds the hypothesis \code{h : False} and then proves the result. Explain precisely what has changed, and why the fact that Lean accepts the new file does not answer the original question.\end{exercise}
\begin{exercise}\label{ex:16.5}Design a review in three steps that detects, in turn, a changed domain, a missing exclusion (such as the condition that a second out-neighbor is not a first out-neighbor), and an unfinished proof containing \code{sorry}. For each step, say what you examine: which part of the Lean file, or which output of Lean.\end{exercise}
\begin{exercise}\label{ex:16.6}Explain the difference between three kinds of check: running a finite test in Python, proving a statement about specific objects in Lean, and proving a universal theorem in Lean. Give an example of each, different from the examples in Section~\ref{ch16:sec:axioms}, and say which cases your example covers.\end{exercise}

